% ===================================================================
%  तालिका सं. 11 — नगर और उसकी इमारतें। — आग।
%  Traduction 2026 d'apres texto/io, controlee sur texto/fr et
%  texto/en. Consigne : voir texto/hi/10-talika-01.tex.
% ===================================================================

%%K t11-apar-1 apar
\VUsaut{2.0mm}
\VUcentre{13.2pt}{90}{तीसरी शृंखला}
\VUsaut{3.0mm}
\VUfilet{16mm}
\VUsaut{5.6mm}
\VUtitre{15.0pt}{60}{तालिका सं. 11}
\VUsaut{1.4mm}
\VUpk{11.4pt}{40}{नगर और उसकी इमारतें। --- आग।}
\VUsaut{2.2mm}

%%K t11-tit-1 sub
\VUpk{10.2pt}{40}{बाहरी बस्ती।}
\VUsaut{3.0mm}

%%K t11-c1-01-1 p
1. --- मैं महासागर से सौ किलोमीटर दूर एक बड़े नगर में रहता हूँ, जो फिर भी
पहले दर्जे का \VUgras{बंदरगाह} \textsuperscript{(1)} है। जो नदी उसके साथ बहती है
वह चार सौ मीटर चौड़ी है और एक बहुत सुंदर लंगरगाह बनाती है।

%%K t11-c1-02-1 p
2. --- \VUgras{घाटों} \textsuperscript{(2)} के साथ का सारा भाग बिलकुल सामुद्रिक और
औद्योगिक रूप रखता है, और एक ऐसी \VUgras{बाहरी बस्ती} \textsuperscript{(3)} बनाता
है जिसमें केवल नाविक और मज़दूर रहते हैं। ये \VUgras{कारख़ानों}
\textsuperscript{(4)} में और \VUgras{गैस-घर} \textsuperscript{(5)} में काम करते हैं,
जिसकी ऊँची \VUgras{चिमनी} \textsuperscript{(6)} और \VUgras{गैस-टंकी} दूर से ही
दिखाई देती है।

%%K t11-c1-03-1 p
3. --- मध्यकाल में नगर किलेबंद था, और उसकी प्राचीरों के कुछ अवशेष अब भी मिलते
हैं, विशेषकर वह \VUgras{द्वार} \textsuperscript{(8)} जो पुराने
\VUgras{गढ़-दुर्ग} \textsuperscript{(7)} का है, और जिसके दोनों ओर उसकी दो \VUgras{मीनारें} \textsuperscript{(9)}
हैं, जो अब \VUgras{कारागार} के काम आती हैं।

%%K t11-c1-04-1 p
4. --- इस सारे \VUgras{मुहल्ले} में, जो बहुत पुराना लगता है, केवल एक इमारत
अपेक्षाकृत नई है: \VUgras{चुंगी-घर} \textsuperscript{(10)}। वह घाट और उस छोटी
\VUgras{गली} \textsuperscript{(11)} के बीच खड़ा है जो पुराने \VUgras{नगर-भवन}
\textsuperscript{(12)} तक जाती है। उस अंतिम इमारत को उस विशाल \VUgras{घंटा-मीनार}
\textsuperscript{(13)} से सहज ही पहचाना जा सकता है जो पुराने नगर पर छाई रहती है।

%%K t11-c1-05-1 p
5. --- गली के दूसरी ओर चुंगी-घर की ही शैली की एक इमारत \VUgras{खड़ी} है: वह
\VUgras{व्यापार-भवन} \textsuperscript{(14)} है। उसका एक मुख एक छोटे चौक की ओर
खुलता है जहाँ \VUgras{प्रशासन-भवन} \textsuperscript{(15)} है। हमारा नगर, वस्तुतः,
राजधानी न होते हुए भी एक महत्त्वपूर्ण ज़िला-मुख्यालय है।

%%K t11-tit-2 sub
\VUsaut{5.0mm}
\VUpk{10.2pt}{40}{नगर का ऊपरी भाग।}
\VUsaut{3.4mm}

%%K t11-c1-06-1 p
6. --- छावनी, जो काफ़ी बड़ी है, विशाल और बहुत स्वास्थ्यकर \VUgras{बैरकों}
\textsuperscript{(16)} में रहती है; वे \VUgras{नगर-उपवन} \textsuperscript{(17)} के
जँगले तक फैली हैं। जो \VUgras{चौड़ी सड़क} \textsuperscript{(19)} इस उपवन तक जाती
है, वह \VUgras{बचत-बैंक} \textsuperscript{(18)} के पीछे से गुज़रकर
\VUgras{महागिरजे} \textsuperscript{(20)} के चौक पर निकलती है।

%%K t11-c1-07-1 p
7. --- ये सब इमारतें एक छोटी पहाड़ी पर सीढ़ीदार उठती हैं, जहाँ हमारा बड़ा
\VUgras{माध्यमिक विद्यालय} \textsuperscript{(21)} भी है।

%%K t11-c1-07-2 p
यदि आप उस हलकी ढाल वाली सड़क से उतरें जो बड़ी सड़क से मिलती है, तो आप
\VUgras{न्यायालय} \textsuperscript{(22)} तक पहुँचते हैं, जिसके सामने एक सुहावना
\VUgras{सार्वजनिक बग़ीचा} \textsuperscript{(26)} है। यह \VUgras{चौक} बहुत बड़ा नहीं,
पर नगर के बीच वह हरियाली और ठंडक का मरूद्यान बनाता है। सो नगरवासी उसमें ख़ूब आते
हैं, जिन्हें \VUgras{कहवाघर} \textsuperscript{(25)} के छज्जे के नीचे बैठकर उस
सैनिक बैंड को सुनना भाता है जो बृहस्पतिवार और रविवार को लॉनों और क्यारियों के
बीच बने \VUgras{बाजा-मंच} \textsuperscript{(27)} में बजता है।

%%K t11-c1-08-1 p
8. --- उन्हीं लॉनों में से एक पर काँसे की एक \VUgras{प्रतिमा}
\textsuperscript{(28)} खड़ी है, हमारे एक नगरवासी की स्मृति में उठाया गया स्मारक,
जिसने अपनी जन्मभूमि की बड़ी सेवा की।

%%K t11-tit-3 sub
\VUsaut{4.0mm}
\VUpk{10.2pt}{40}{आना-जाना।}
\VUsaut{3.4mm}

%%K t11-c1-09-1 p
9. --- हमारा नगर अभी बिजली से रोशन नहीं; उसमें केवल \VUgras{गैस के लैंप}
\textsuperscript{(29)} हैं। पर \VUgras{स्थानीय अख़बार} इस बात के लिए मुहिम चला रहे
हैं कि रोशनी की वह व्यवस्था सुधारी जाए, जैसे ट्रामों की
\VUgras{खिंचाई की व्यवस्था} पहले ही सुधारी जा चुकी है। \VUgras{घोड़ों} से
\VUgras{खींची} जाने वाली पुरानी गाड़ियों की जगह अब व्यावहारिक
\VUgras{बिजली की ट्रामें} \textsuperscript{(31)} आ गई हैं, जो बहुत \VUgras{कम}
किराए में यात्रियों को \VUgras{नगर के} एक सिरे से \VUgras{दूसरे तक} जल्दी पहुँचा
देती हैं।

%%K t11-c1-10-1 p
10. --- न्यायालय के पास \VUgras{बैंक} \textsuperscript{(32)} है, जहाँ व्यापारिक
हुंडियाँ भुनाई जाती हैं। \VUgras{बैंक के हरकारे} \textsuperscript{(33)} लोगों के
घर-घर जाकर बकाया वसूलते हैं।

%%K t11-tit-4 sub
\VUsaut{4.0mm}
\VUpk{10.2pt}{40}{कला और विद्या।}
\VUsaut{3.4mm}

%%K t11-c1-11-1 p
11. --- पास की वह सुंदर इमारत \VUgras{संग्रहालय} है। उसमें एक साथ नगर का
\VUgras{पुस्तकालय} \textsuperscript{(34)}, कुछ सुंदर
\VUgras{प्राकृतिक इतिहास के संग्रह} \textsuperscript{(35)} (एक प्राकृतिक इतिहास
संग्रहालय) और एक सुडौल \VUgras{चित्रशाला} \textsuperscript{(36)} है, जो एक शानदार
स्थायी \VUgras{प्रदर्शनी} बनाती है।

%%K t11-c1-11-2 p
वह सप्ताह में दो बार जनता के लिए खुलती है, पर दर्शक किसी भी दिन उसे देख सकते हैं,
\VUgras{रक्षक} \textsuperscript{(37)} को थोड़ी बख़्शीश देकर। \VUgras{चित्रकार}
\textsuperscript{(38)} वहाँ काम करने आते हैं। उन्हें अपनी तिपाइयों के
\VUgras{सामने} देखा जा सकता है, हाथ में \VUgras{रंग-पट्टी} लिए, प्रसिद्ध चित्रों
की नक़ल उतारते।

%%K t11-c1-12-1 p
12. --- संग्रहालय के पीछे की ऊँचाइयों पर वह \VUgras{गुंबद}
\textsuperscript{(40)} जो \VUgras{चिकित्सालय} \textsuperscript{(39)} का है, और उस
\VUgras{अध्यापक-प्रशिक्षण विद्यालय} \textsuperscript{(41)} की इमारतें मुश्किल से
दिखाई देती हैं जहाँ हमारे नगर-विद्यालयों के अध्यापक प्रशिक्षित हुए थे। रंगमंच के
चौक पर एक \VUgras{डाकघर} \textsuperscript{(43)} है, जो एक \VUgras{निजी मकान}
\textsuperscript{(42)} में बैठा है।

%%K t11-tit-5 sub
\VUsaut{5.0mm}
\VUpk{11.4pt}{40}{आग।}
\VUsaut{3.0mm}

%%K t11-tit-6 sub
\VUpk{10.2pt}{40}{आग की हाँक।}
\VUsaut{3.4mm}

%%K t11-c2-01-1 p
1. --- नगर में डरावनी ख़बर फैल रही है: रंगमंच में अभी-अभी आग लग गई! जब तक मैं
\VUgras{चौक} \textsuperscript{(96)} पहुँचता हूँ, भीड़ \VUgras{पटरी}
\textsuperscript{(46)} और \VUgras{सड़क} \textsuperscript{(47)} पर जुट चुकी है।
\VUgras{डाक के बाबू} \textsuperscript{(44)} और \VUgras{डाकिए} \textsuperscript{(45)}
डाकघर से बाहर आ रहे हैं। एक \VUgras{ट्राम-चालक} \textsuperscript{(48)} को अपनी
ट्राम रोकनी पड़ी है ताकि नगर की एक \VUgras{रोगी-गाड़ी} \textsuperscript{(50)}
निकल सके, जो एक \VUgras{शल्य-चिकित्सक} \textsuperscript{(52)} और एक
\VUgras{परिचारिका} \textsuperscript{(51)} को लेकर आ रही है, एक \VUgras{घायल}
\textsuperscript{(53)} की देखभाल के लिए, जिसे \VUgras{स्ट्रेचर} \textsuperscript{(54)}
पर \VUgras{अख़बार के खोखे} \textsuperscript{(55)} के पास लाया गया है।

%%K t11-c2-02-1 p
2. --- कवायद से लौटती पैदल सेना की एक टुकड़ी रुक गई है, ताकि घोड़ों पर सवार
\VUgras{नगर-रक्षकों} \textsuperscript{(61)} के साथ व्यवस्था बनाए रखने में मदद करे।
\VUgras{घुड़सवारों} के घोड़े दो पैरों पर खड़े हो जाते हैं; वे \VUgras{सैनिकों}
\textsuperscript{(57)} या \VUgras{जेंडार्मों} \textsuperscript{(63)} से बेहतर
\VUgras{तमाशबीनों} \textsuperscript{(56)} को पीछे धकेल पाते हैं। जेंडार्म वैसे भी
बहुत व्यस्त हैं। उनमें से एक \VUgras{पुलिस के सिपाहियों} \textsuperscript{(64)} को
बता रहा है कि एक और \VUgras{घायल} \textsuperscript{(65)} को कहाँ ले जाना है, एक
वीर दमकल वाला जो सीढ़ी से गिर पड़ा है।

%%K t11-c2-03-1 p
3. --- \VUgras{अफ़सर} \textsuperscript{(66)} ठीक-ठीक बताता है कि जो
\VUgras{भाप-पंप} \textsuperscript{(67)} अभी आ रहा है उसे कहाँ रखना है। उस
शक्तिशाली यंत्र की प्रतीक्षा में उसने एक \VUgras{हाथ-पंप} \textsuperscript{(68)}
लगवा दिया है, जिसकी लंबी \VUgras{नली} \textsuperscript{(69)} आग पर पानी की धाराएँ
उड़ेल रही है। छह दमकल वाले उसे चला रहे हैं, जब तक उनका साथी \VUgras{टोंटी}
\textsuperscript{(70)} थामे \VUgras{धार} \textsuperscript{(71)} को आग के बीच साधता
है।

%%K t11-c2-04-1 p
4. --- रंगमंच के दूसरी ओर बड़ी \VUgras{सीढ़ी} \textsuperscript{(72)} खड़ी की गई
है। वह दमकल वालों को उन लपटों के पास पहुँचने देती है जो काले \VUgras{धुएँ}
\textsuperscript{(75)} के बगूलों के बीच उठती हैं।

%%K t11-tit-7 sub
\VUsaut{5.0mm}
\VUpk{10.2pt}{40}{वीरता के काम।}
\VUsaut{3.4mm}

%%K t11-c2-05-1 p
5. --- \VUgras{आग} \textsuperscript{(74)} \VUgras{रंगमंच} \textsuperscript{(73)} के
प्रवेश-कक्ष में लगी, उस \VUgras{ओपेरा} के पूर्वाभ्यास के समय जिसका पहला
\VUgras{प्रदर्शन} कल होने वाला था। सौभाग्य से सभा-कक्ष में केवल कुछ
\VUgras{दर्शक} \textsuperscript{(76)} थे। बेचारों ने \VUgras{बालाख़ाने}
\textsuperscript{(77)} में शरण ली, जहाँ वीर \VUgras{दमकल वाले} \textsuperscript{(86)}
सीढ़ी और रस्सियों के साथ उनकी मदद को आ रहे हैं।

%%K t11-c2-06-1 p
6. --- \VUgras{द्वार-मंडप} \textsuperscript{(78)} के नीचे, जहाँ
\VUgras{नाट्य-सूचना} \textsuperscript{(79)} प्रदर्शन की घोषणा करती है,
\VUgras{बाजेवाले} \textsuperscript{(81)} भागते दिखाई देते हैं, जो
\VUgras{वाद्य-वृंद} \textsuperscript{(80)} के हैं, अपने वाद्य उठाए:
\VUgras{वायलिन} \textsuperscript{(81)}, \VUgras{बाँसुरी} \textsuperscript{(82)},
\VUgras{चेलो} \textsuperscript{(83)}, \VUgras{ट्रॉम्बोन} \textsuperscript{(84)},
\VUgras{ढोल} \textsuperscript{(85)}, इत्यादि। \VUgras{साज-सामान}
\textsuperscript{(87)} का एक हिस्सा बचाने का यत्न किया जा रहा है।

%%K t11-c2-07-1 p
7. --- \VUgras{अभिनेताओं} \textsuperscript{(89)} और \VUgras{अभिनेत्रियों}
\textsuperscript{(88)} को, गायकों और गायिकाओं को, अपने मंच के \VUgras{वस्त्रों}
\textsuperscript{(90)} में ही भागना पड़ा है। मुख्य गायक एक \VUgras{पत्रकार}
\textsuperscript{(92)} को समझा रहा है कि यह कैसे हुआ। \VUgras{संवाददाता} पहले ही
क्षण से अधिकारियों के साथ दौड़ा आया था: \VUgras{प्रशासक} \textsuperscript{(91)},
\VUgras{नगराध्यक्ष} \textsuperscript{(93)}, \VUgras{पुलिस अधीक्षक}
\textsuperscript{(94)} और एक \VUgras{न्यायाधीश} \textsuperscript{(95)}, जो जाँच
करेंगे कि ज़िम्मेदारी किसकी है।
\VUsaut{6.0mm}
\VUfilet{22mm}
